\documentclass[10pt,a4paper,sans]{moderncv}

\moderncvstyle{classic}
\moderncvcolor{blue}

\usepackage[scale=0.82]{geometry}
\setlength{\hintscolumnwidth}{1.8cm}
\usepackage{setspace}
\setstretch{0.78}

\name{Patrick Serrano}{}
\title{Engenheiro de Software Sênior Backend (Go · Clojure / Sistemas Financeiros)}
\email{patrickcmserrano@gmail.com}
\social[linkedin]{patrickcmserrano}
\social[github]{patrickcmserrano}

\begin{document}
\makecvtitle

\section{Resumo}
Interrompi o curso de Física na UFF para me tornar desenvolvedor --- não por acidente, mas porque reconheci que construir sistemas com rigor era o que queria fazer de verdade. Passei os últimos 7 anos numa consultoria construindo um motor de pagamentos multi-adquirente do zero, da modelagem de dados à integração com 25+ gateways financeiros --- sistemas críticos em produção onde disponibilidade, resiliência e rastreabilidade não são opcionais. Hoje estou construindo em Go minha própria plataforma de gestão de patrimônio, o que me fez querer trabalhar em produto onde as decisões de arquitetura têm consequências que você acompanha no dia seguinte. Trabalho melhor em times onde código é discutido antes de ser entregue: testes como documentação executável, code reviews como design colaborativo, abstrações que emergem do problema em vez de preceder ele.

\section{Habilidades}
\cvitem{Ecossistema Clojure}{Pathom 3, Datomic (Datalog), Clara Rules, Malli, Reitit, Pedestal, Onyx, Fulcro, Fulcro RAD, Reagent, shadow-cljs}
\cvitem{Linguagens}{Clojure, ClojureScript, Go, TypeScript, Svelte, Python}
\cvitem{Banco de Dados}{PostgreSQL, Datomic, Redis}
\cvitem{Arquitetura}{Microservices, Event-Driven, CQRS, Clean Architecture, Polylith, ETL Pipelines}
\cvitem{Integrações}{Gateways: Adyen, Cielo, eRede, Pagar.me, Mercado Pago, Getnet, Tuna, Unico; Antifraude: ClearSale, Konduto; E-commerce: VTEX, Loja Integrada}
\cvitem{Mensageria}{NATS JetStream, Onyx; Kafka (familiaridade teórica, experiência com sistemas análogos)}
\cvitem{Infra / DevOps}{AWS (SQS, Lambda, ECS, SSM), Docker, Pulumi, CI/CD, Github Actions}

\section{Experiência Profissional}

\cventry{2025 -- Atual}{Plataforma de Gestão de Patrimônio em Cripto (Ark Streams)}{Projeto Independente}{}{}{%
Sistema financeiro em Go para coleta, processamento e execução de ordens em múltiplas exchanges, com requisitos não-funcionais tratados como cidadãos de primeira classe: disponibilidade contínua, resiliência a falhas de rede, latência controlada e rastreabilidade de estado. Produto com visão além do uso pessoal --- democratizar para pessoas físicas o que hoje só instituições financeiras operam.
\begin{itemize}
    \item \textbf{Disponibilidade e resiliência por design:} coletor headless rodando 24/7 em VPS, checkpoint periódico do estado analítico para recuperação instantânea no restart, reconexão transparente após queda de rede --- sistema financeiro não pode falhar no meio de uma operação, então resiliência foi projetada desde o início, não adicionada depois;
    \item \textbf{Escalabilidade e independência de plataforma:} pipeline multi-exchange (Bitget, Binance, Bybit, OKX + Yahoo Finance) com eventos normalizados no barramento NATS JetStream --- adicionar ou trocar uma fonte não afeta nada downstream. Backtest e produção executam o mesmo código sobre os mesmos eventos, eliminando divergência por propriedade estrutural;
    \item \textbf{Performance medida e documentada:} benchmarks de pipeline ($\sim$66\,µs / 100 candles) e de transporte (migração HTTP $\rightarrow$ IPC nativo: 70--78\% menos RAM, 100$\times$ menos latência de streaming). Cada decisão de otimização precedida por medição, não por intuição;
    \item \textbf{Integrabilidade:} API REST (Chi v5), WebSocket para streaming de eventos, PostgreSQL para persistência relacional --- arquitetura pensada para integrar com ecossistemas externos, não operar em silos;
    \item \textbf{Do sinal à execução no mesmo sistema:} order execution integrado com Bitget --- ordens reais, gestão de SL/TP, controles de alavancagem e margem. Inteligência e execução no mesmo ambiente, sem fricção entre análise e ação.
\end{itemize}
Tecnologias: Go 1.23, PostgreSQL, NATS JetStream, Chi v5, WebSocket, Wails v2, Svelte 5, Docker, GitHub Actions
}

\cventry{2022 -- Jul 2025}{Motor de Pagamentos (Octopus Pay)}{Consultoria em Software}{}{}{%
Plataforma multi-adquirente de alta disponibilidade para orquestração de pagamentos com 25+ integrações ativas.\newline{}
Responsabilidades
\begin{itemize}%
    \item Design e implementação de 25+ adaptadores para gateways (Adyen, Cielo, eRede, Pagar.me, Mercado Pago, Getnet, Tuna, Unico), plataformas (VTEX, Loja Integrada) e antifraude (ClearSale, Konduto) via modelo canônico EQL/Pathom;
    \item Desenvolvimento de motor ETL para ingestão e processamento de milhões de eventos financeiros com orquestração extensível (Onyx);
    \item Liderança da migração Pathom 2 → Pathom 3 e derivação de atributos calculados em tempo real com EQL, com disseminação das práticas para o time via code reviews e pair programming;
    \item Desenvolvimento de APIs REST com Reitit + Pedestal, validação Malli, Swagger/OpenAPI e padrões de resiliência (circuit breaker);
    \item Desenvolvimento e manutenção de suíte de testes com \texttt{clojure.test} e Midje (BDD): testes unitários de modelos e resolvers Pathom, testes de integração com servidor Pedestal real (fixtures de lifecycle, mocks via \texttt{clj-http.fake}), testes de contratos de API por rota e validação Malli de schemas --- disciplina de criar teste para cada bug corrigido;
    \item Desenvolvimento de dashboards administrativos full-stack em ClojureScript + Fulcro RAD com integração Datomic/Datalog;
    \item Modelagem de regras de forward-chaining com Clara Rules para cálculos financeiros e workflows automatizados;
    \item Gestão de infraestrutura AWS via Pulumi (infra-as-code), configuração via SSM e pipelines de deploy automatizados.
\end{itemize}
Tecnologias: Clojure, ClojureScript, Datomic, Pathom 3, Clara Rules, Fulcro RAD, Reitit, Pedestal, AWS, Pulumi
}

\cventry{2020 -- 2022}{Marketplace B2B (Zougue MPMS)}{Consultoria em Software}{}{}{%
Gestão de marketplace integrada com VTEX para sellers e operadores.\newline{}
Responsabilidades
\begin{itemize}%
    \item Desenvolvimento de endpoints de integração VTEX: simulação de carrinho, SKU, preço/estoque, notas fiscais, autorização de despacho;
    \item Modelagem de esquemas Datomic para domínios centrais (order, seller, doca, estoque, contrato) com Pathom resolvers multi-seller;
    \item Implementação de regras Clara Rules para workflows, roteamento e cálculos de faturamento/price-tags;
    \item Desenvolvimento de sistema de importação de catálogo via XLSX com validação Malli e relatórios de erro;
    \item Desenvolvimento de portal administrativo em ClojureScript/Fulcro com filtros de workflow e listagem de pedidos;
    \item Desenvolvimento de portal Marketplace multi-tenant com white-label, autenticação magic link e Metabase integrado.
\end{itemize}
Tecnologias: Clojure, ClojureScript, Fulcro, Pathom, Datomic, Clara Rules, VTEX, Metabase
}

\cventry{2018 -- 2019}{Projetos Iniciais (Plataforma Imobiliária · Gateway de Pagamentos)}{Consultoria em Software}{}{}{%
\begin{itemize}
    \item Desenvolvimento fullstack de aplicação mobile e web do zero em ClojureScript/Fulcro com React Native (plataforma imobiliária);
    \item Desenvolvimento de front-end e integrações com sistemas financeiros em ClojureScript/Fulcro para gateway de pagamentos.
\end{itemize}
Tecnologias: ClojureScript, Fulcro, React Native, REST APIs
}

\section{Formação Acadêmica}
\cventry{2024 -- Atual}{Bacharelado em Engenharia de Software}{Descomplica Faculdade Digital}{}{}{}
\cventry{2015 -- 2018}{Bacharelado em Física}{Universidade Federal Fluminense}{}{}{}

\section{Formação Complementar}
\cvitem{Full Cycle 3.0}{Segurança \& Identidade: OAuth 2.0, OpenID Connect, Keycloak, JWT · Mensageria: EDA, Kafka · Arquitetura: Microsserviços, Hexagonal, Clean Architecture, DDD, SOLID · DevOps: Kubernetes, Terraform, Observabilidade, CI/CD}

\section{Idiomas}
\cvitemwithcomment{Português}{Nativo}{}
\cvitemwithcomment{Inglês}{Avançado}{}
\cvitemwithcomment{Alemão}{Básico}{}
\cvitemwithcomment{Espanhol}{Básico}{}

\end{document}
